โครงงานนี้นำเสนอการออกแบบและพัฒนาระบบจัดการสินค้าคงคลังสำหรับธุรกิจซื้อมา-ขายไปขนาดเล็กและขนาดกลาง ระบบนี้รองรับกระบวนการทำงานของธุรกิจตัวกลางที่ซื้อสินค้าจากผู้จัดจำหน่าย รับสินค้าเข้าคลัง ขายสินค้าให้ลูกค้า และติดตามเอกสารทางธุรกิจที่เกี่ยวข้องกับลูกหนี้และเจ้าหนี้

ระบบถูกพัฒนาในรูปแบบเว็บแอปพลิเคชันสามชั้น โดยใช้ React และ Vite สำหรับส่วนติดต่อผู้ใช้ ใช้ Django และ Django REST Framework สำหรับระบบหลังบ้าน และใช้ MySQL สำหรับการจัดเก็บข้อมูลแบบเชิงสัมพันธ์ ฟังก์ชันที่พัฒนาแล้วประกอบด้วยการจัดการข้อมูลหลักของสินค้า หมวดหมู่ ผู้จัดจำหน่าย และลูกค้า การจัดการรายการซื้อและขาย ใบเสนอราคา ใบวางบิล ชุดการชำระเงิน ใบลดหนี้ การตรวจสอบสต็อก แดชบอร์ดสรุปข้อมูล และผู้ช่วยแบบอ่านข้อความสำหรับคำถามเชิงปฏิบัติการบางประเภท

ตรรกะทางธุรกิจที่สำคัญถูกตรวจสอบในฝั่งระบบหลังบ้าน ปริมาณสินค้าคงคลังคำนวณจากจำนวนสินค้าที่รับเข้าจากรายการซื้อ หักด้วยจำนวนสินค้าที่ถูกผูกไว้กับรายการขาย แทนการบันทึกสต็อกเป็นค่าคงที่ในข้อมูลสินค้า ระบบตรวจสอบสต็อกในฝั่งเซิร์ฟเวอร์ก่อนให้รายการขายเปลี่ยนเป็นสถานะที่ตัดสต็อก และจัดสรรต้นทุนจากชั้นสินค้าที่รับเข้าด้วยหลัก FIFO นอกจากนี้ ระบบยังใช้จุดเชื่อมต่อสำหรับตรวจสอบสิทธิ์ของรายการที่สามารถนำไปสร้างใบวางบิล ชุดการชำระเงิน และใบลดหนี้ได้

ผลลัพธ์ของโครงงานคือระบบจัดการสินค้าคงคลังบนเว็บที่สามารถเชื่อมโยงงานจัดซื้อ งานขาย การมองเห็นสต็อก และการติดตามเอกสารทางการเงินไว้ในระบบเดียว ความสามารถของผู้ช่วยในปัจจุบันจำกัดอยู่ที่คำถามเกี่ยวกับสต็อกและการส่งมอบ สรุปลูกค้าหรือผู้จัดจำหน่าย ตรวจสอบลูกหนี้ เจ้าหนี้ และรายการผิดปกติ รวมถึงค้นหาเอกสารหรือรายละเอียดรายการสินค้า ระบบนี้ช่วยสนับสนุนการทำงานประจำวันของธุรกิจซื้อมา-ขายไป และสามารถพัฒนาต่อยอดได้ในอนาคต เช่น การเพิ่มสูตรวางแผนสินค้าคงคลัง การกำหนดสิทธิ์ตามบทบาทผู้ใช้ บันทึกประวัติการเปลี่ยนแปลงอย่างเป็นทางการ และการทดสอบในระดับการใช้งานจริงที่ใหญ่ขึ้น
